\lecture{14}{Thu 21 Nov 2024 09:31}{Ideals and Quotient Rings}
An ideal of a ring is analagous to a normal subgroup of a group.
\begin{definition}[Ideal]\label{dfn:IDL}
	An ideal $I$ in a ring $R$ is a subring $I\subseteq R$ such that if $a\in I$ and $r\in R$, $ar,ra\in I$. That is, $rI\subseteq I$ and $Ir\subseteq I$ for all $r\in R$.
\end{definition}
Every ring $R$ has at least two ideals, $R$ and $\left\{ 0 \right\} $. These are the trivial ideals. If $R$ is a ring with unity, and $I\lhd  R$ is an ideal with the unity of $R$ contained in $I$, then $I=R$ because $r1\in I$ for all $r$ by definition of an ideal.
\begin{definition}\label{dfn:kjs}
	Let $R$ be a commutative ring with unity, and let $\left\langle a \right\rangle =\left\{ ar \mid r\in R \right\} $.
\end{definition}
\begin{theorem}\label{thm:kjsd}
	Let $a\in R$ be a commutative ring with unity.  Then $\left\langle a \right\rangle $ is an ideal in $R$.
\end{theorem}
\begin{proof}
	Note that $a1=a\in \left\langle a \right\rangle $. Note that $ar_1-ar_2=a\left( r_1-r_2) \right) \in \left\langle a \right\rangle $. Finally, note $r_2ar_1=ar_1r_2\in \left\langle a \right\rangle $. Finally, $ar_1ar_2=a\left( r_1ar_2 \right) \in \left\langle a \right\rangle $. This is an ideal
\end{proof}
\begin{definition}[Principal]\label{dfn:PI}
	An ideal $I$ is \textbf{principal} if $I=\left\langle a \right\rangle $ for some $a\in R$.
\end{definition}

For example, every ideal in $\mathbb{Z}$ is principal. If $I$ is nontrivial, $I$ must contain some $m\in\mathbb{Z}^+$. Then there must be a least positive integer $n>0\in I$. Let $a\in I$. We have $a=nq+r$ for some $q,r\in\mathbb{Z}$ and $0\leq r<n$. Then $r=a-nq$. Since $a,nq\in I$, $r\in I$. Therefore, $r=0$, and $a=nq\in \left\langle n \right\rangle $. Therefore, $\left\langle n \right\rangle =I$.

Note that not all ideals are principal. Consider $I=\left\{ 2a+xb \mid a,b\in\mathbb{Z} \right\} $, an ideal of $\mathbb{Z}[x]$. This ideal is generated by two elements.
\begin{theorem}\label{thm:bwaaa}
	Let $I$ be an ideal of $R$. The quotient group $R /I$ is a ring with multiplication defined by $\left( r+I \right) \left( s+I \right) =rs+I$.
\end{theorem}
\begin{proof}
	$R /I$ is a quotient group because $(R,+)$ is abelian. We need to prove it now forms a ring. Let $r+I,s+I,t+I\in R /I$. First, we show multiplication is well defined. Let $r+I=s+I$. Since $r'\in r+I$ by properties of cosets, there exists $a\in I$ such that $r'=r+a$. Similarly, there exists $b\in I$ such that $s'=s+b$. We now have
	\[
		\left( r'+I \right) \left( s'+I \right) =r's'+I=\left( r+a \right) \left( s+b \right) +I=(rs+as+rb+ab)+I
	.\]
	We want to show $rs+as+rb+ab\in I$. Recall that $a,b\in I$, so any product with them is in $I$. Therefore, $rs,as,rb,ab\in I$, and since $I$ is a ring, $rs+as+rb+ab\in I$. Therefore, $r's'+I =rs+I$.

	Closure is trivial. Consider associativity.
	\[
		\left( r+I \right) \left( \left( s+I \right) \left( t+I \right)  \right) =\left( r+I \right) \left( st+I \right) =r(st)+I=(rs)t+I=\left( \left( r+I \right) \left( s+I \right)  \right) \left( t+I \right) 
	.\]
	Therefore, $R /I$ is associative. Now consider distributivity.
	\[
		\left( r+I \right) \left( s+I+t+I \right) =\left( r+I \right) \left( (s+t)+I \right) =r(s+t)+I=\left( rs+rt \right) +I=\left( rs+I \right) +\left( rt+I \right) 
	.\]
	The proof is the same on the other side. Therefore, $R /I$ is a ring.
\end{proof}
\begin{definition}[Quotient Ring]\label{dfn:QR}
	Let $R$ be a ring and let $I\subseteq R$ be an ideal. Then $R  /I$ is a ring with addition defined as
	\[
		\left( r+I \right) +\left( s+I \right) =\left( r+s \right) +I
	,\]
	and multiplication defined as
	\[
		\left( r+I \right) \left( s+I \right) =rs+I
	.\]
\end{definition}

For example, somehow $\mathbb{Z}_3[x] / \left\langle x^2+1 \right\rangle $ is a quotient ring. The identity is $\left\langle x^2+1 \right\rangle $. The multilpicative identity is $1+\left\langle x^2+1 \right\rangle $, and the other elements are of the form $p(x)+\left\langle x^2+1 \right\rangle  $. Also, consider
\[
	x^2+\left\langle x^2+1 \right\rangle =2+\left\langle x^2+1 \right\rangle 
.\]
Since $x^2+1\in \left\langle x^2+1 \right\rangle $ is the identity, $x^2-\left( x^2+1 \right) +\left\langle x^2+1 \right\rangle =-1+\left\langle x^2+1 \right\rangle $. Then $-1\equiv 2\mod{3}$.
